%% toc: P3 40 40.4
\section{Proofs for constant-probability regimes and their countable models}
\label{sec:annex-p03-w03-p03-067-constant-p}

Throughout, the assumptions of \cref{def:p03-w03-p03-067-assumptions} are in
force.  Every language used is finite, so compactness
(\cref{thm:fol-compactness}) and completeness (\cref{thm:fol-completeness})
apply without a choice principle; the downward L\"owenheim--Skolem theorem
(\cref{thm:downward-ls}) is used with the Axiom of Choice.  For a graph
\(G\) and vertices \(a,b\), \(a\not\sim b\) abbreviates \(\neg\,a\sim b\).

\begin{annexitem}{annex:p03-w03-p03-067-as-theory}{thm:p03-w03-p03-067-as-theory}{The almost sure theory of a zero--one law}
\begin{proof}
\emph{Finite conjunctions.}  Let \(A_1,\ldots,A_m\in T_{\mathrm{as}}\).  For
each \(n\),
\(\Pr[\neg(A_1\land\cdots\land A_m)]\leq\sum_{i\leq m}\Pr[\neg A_i]\), and each
summand tends to \(0\).  So the conjunction holds almost surely.  In
particular its probability is positive for some \(n\), and then some
structure of size \(n\) satisfies every \(A_i\).

\emph{Item 1.}  By the previous paragraph every finite subset of
\(T_{\mathrm{as}}\) has a model, so \(T_{\mathrm{as}}\) has a model by clause~1
of \cref{thm:fol-compactness}; it is consistent.  Suppose
\(T_{\mathrm{as}}\entails B\).  By clause~2 of \cref{thm:fol-compactness}
there are \(A_1,\ldots,A_m\in T_{\mathrm{as}}\) with
\(A_1,\ldots,A_m\entails B\).  Every structure satisfying the conjunction
satisfies \(B\), so \(\Pr[B]\geq\Pr[A_1\land\cdots\land A_m]\to1\) and
\(B\in T_{\mathrm{as}}\).  For completeness, let \(B\) be any sentence.  By
hypothesis \(B\) holds almost surely, so \(B\in T_{\mathrm{as}}\), or almost
never, so \(\neg B\) holds almost surely and \(\neg B\in T_{\mathrm{as}}\).
Both cannot occur, since \(T_{\mathrm{as}}\) is consistent.

\emph{Item 2.}  For \(r\geq1\) let \(D_r\) say that there exist pairwise
distinct \(x_1,\ldots,x_r\).  Every structure with exactly \(n\geq r\)
elements satisfies \(D_r\), so \(\Pr[D_r]=1\) for all \(n\geq r\) and
\(D_r\in T_{\mathrm{as}}\).  A model of \(T_{\mathrm{as}}\) therefore has at
least \(r\) elements for every \(r\), and is infinite.

\emph{Item 3.}  Let \(\mathcal M\) be a model of \(T_{\mathrm{as}}\), which
exists by item~1.  The language is finite, so \cref{thm:downward-ls} with
\(A=\varnothing\) gives an elementary substructure
\(\mathcal N\preccurlyeq\mathcal M\) of at most countable size.  Elementary
substructures satisfy the same sentences, so \(\mathcal N\) is a model of
\(T_{\mathrm{as}}\), hence infinite by item~2, hence countably infinite.
\end{proof}
\end{annexitem}

\begin{annexitem}{annex:p03-w03-p03-067-witness}{thm:p03-w03-p03-067-witness}{Every extension axiom holds almost surely}
\begin{proof}
Fix \(n\geq r+s\) and the vertex set \(V\) of \(G(n,p)\).  Whether a choice
of distinct \(x_1,\ldots,x_r,y_1,\ldots,y_s\) has a witness depends only on
the disjoint sets \(X=\{x_1,\ldots,x_r\}\) and \(Y=\{y_1,\ldots,y_s\}\).  Fix
such \(X\) and \(Y\).  For \(z\in V\setminus(X\cup Y)\) let \(W_z\) be the event
that \(z\sim x\) for every \(x\in X\) and \(z\not\sim y\) for every \(y\in Y\).
The event \(W_z\) is determined by the \(r+s\) distinct pairs \(\{z,x\}\) and
\(\{z,y\}\); by the independence in \cref{def:p03-w03-p03-067-models},
\(\Pr[W_z]=p^r(1-p)^s=\epsilon\).  For distinct candidates \(z\neq z'\) these
families of pairs are disjoint, because each pair in the family of \(z\)
contains \(z\) and one point of \(X\cup Y\), and the pair \(\{z,z'\}\) belongs
to neither family.  Events determined by disjoint subfamilies of mutually
independent edge events are mutually independent, so the \(n-r-s\) events
\(W_z\) are mutually independent and
\[
 \Pr\Bigl[\,\bigcap_{z\notin X\cup Y}\neg W_z\Bigr]=(1-\epsilon)^{n-r-s}.
\]
The outcome fails \(A_{r,s}\) exactly when some pair \((X,Y)\) has no witness.
There are \(\binom nr\binom{n-r}s\leq n^rn^s\) such pairs, and the union bound
gives \(\Pr[G(n,p)\not\entails A_{r,s}]\leq n^{r+s}(1-\epsilon)^{n-r-s}\).

For the limit, put \(k=r+s\).  If \(k=0\) then \(\epsilon=1\) and the bound is
\(0^n=0\) for \(n\geq1\).  If \(k\geq1\) then \(0<\epsilon<1\) because
\(0<p<1\); put \(q=1-\epsilon\in(0,1)\).  Then
\[
 \ln\bigl(n^kq^{n-k}\bigr)
 =n\Bigl(\frac{k\ln n}{n}+\Bigl(1-\frac kn\Bigr)\ln q\Bigr),
\]
where \(k\ln n/n\to0\) and \((1-k/n)\ln q\to\ln q<0\).  The bracket is
eventually below \(\tfrac12\ln q\), so the logarithm tends to \(-\infty\) and
the bound tends to \(0\).  The remaining axioms of \(T\), irreflexivity and
symmetry, hold in every outcome of \(G(n,p)\).  So every sentence of \(T\)
holds almost surely.
\end{proof}
\end{annexitem}

\begin{annexitem}{annex:p03-w03-p03-067-complete}{thm:p03-w03-p03-067-complete}{Duplicator stays alive; \(T\) is complete}
\begin{proof}
\emph{The strategy.}  Let \(G\) and \(G'\) be models of \(T\) and \(t\geq1\).
Write \((a_k,b_k)\) for the pair of vertices chosen in round \(k\), with
\(a_k\in G\) and \(b_k\in G'\).  Duplicator keeps the following invariant
after every round: for all earlier rounds \(k,l\), \(a_k=a_l\) iff
\(b_k=b_l\), and \(a_k\sim a_l\) iff \(b_k\sim b_l\).  It holds vacuously
before the first round.  By the first clause, \(a_k\mapsto b_k\) is a
well-defined bijection \(\pi\) from the chosen vertices of \(G\) onto those of
\(G'\).

Suppose Spoiler picks \(a\in G\); the case of a vertex of \(G'\) is the same
with the roles exchanged, since both graphs satisfy \(T\).  If \(a=a_k\) for
some earlier \(k\), Duplicator answers \(b_k\), and the invariant is
unchanged because the new pair repeats an old one.  Otherwise \(a\) is new.
Let \(P\) be the set of \(\pi(a_k)\) with \(a_k\sim a\) and \(Q\) the set of
\(\pi(a_k)\) with \(a_k\not\sim a\).  Since \(\pi\) is a bijection, \(P\) and
\(Q\) are disjoint and their union is the set of all earlier answers in
\(G'\).  Let \(r=\abs P\), \(s=\abs Q\), list \(P\) and \(Q\) without
repetition, and apply \(A_{r,s}\) in \(G'\): there is \(b\notin P\cup Q\)
adjacent to every vertex of \(P\) and to none of \(Q\).  Duplicator answers
\(b\).  Then \(a\) and \(b\) are both new, so equality is matched; \(a\) and
\(b\) have the same adjacency to each earlier pair by the choice of \(b\); and
neither is adjacent to itself by irreflexivity.  By symmetry of \(\sim\) in
both graphs, the order of arguments does not matter.  The invariant
therefore holds after every round, and after \(t\) rounds the pairing is a
partial isomorphism: Duplicator wins \(\mathrm{EHR}(G,G';t)\).

\emph{Consistency.}  Let \(\Sigma\) be a finite subset of \(T\).  By
\cref{thm:p03-w03-p03-067-witness} with \(p=1/2\), each member of \(\Sigma\)
fails in \(G(n,1/2)\) with probability tending to \(0\), so for large \(n\)
the probability that all of \(\Sigma\) holds is positive, and some finite
graph satisfies \(\Sigma\).  By clause~1 of \cref{thm:fol-compactness},
\(T\) has a model.

\emph{Completeness.}  Let \(B\) be a sentence of quantifier rank \(t\) (if
\(t=0\), read \(t=1\); agreement through rank \(1\) includes rank \(0\)).  For
any models \(G,G'\) of \(T\), Duplicator wins \(\mathrm{EHR}(G,G';t)\), so by
\cref{thm:p03-ef} the two graphs agree on \(B\).  Fix one model \(G_0\).  If
\(G_0\entails B\) then every model of \(T\) satisfies \(B\), so
\(T\entails B\) and \(T\vdash B\) by \cref{thm:fol-completeness}; otherwise
every model satisfies \(\neg B\) and \(T\vdash\neg B\).
\end{proof}
\end{annexitem}

\begin{annexitem}{annex:p03-w03-p03-067-zero-one}{thm:p03-w03-p03-067-zero-one}{Zero--one law for constant \(p\)}
\begin{proof}
By \cref{thm:p03-w03-p03-067-complete}, \(T\vdash B\) or \(T\vdash\neg B\),
and not both since \(T\) is consistent.  Suppose \(T\vdash B\).  By
soundness and clause~2 of \cref{thm:fol-compactness}, finitely many axioms
\(C_1,\ldots,C_m\) of \(T\) entail \(B\).  Every outcome of \(G(n,p)\) that
fails \(B\) fails some \(C_i\), so
\(\Pr[G(n,p)\not\entails B]\leq\sum_{i\leq m}\Pr[G(n,p)\not\entails C_i]\),
which tends to \(0\) by \cref{thm:p03-w03-p03-067-witness}.  The limit is
\(1\).  If \(T\vdash\neg B\), the same argument applied to \(\neg B\) gives
\(\Pr[G(n,p)\entails\neg B]\to1\), so the limit for \(B\) is \(0\).  Whether
\(T\vdash B\) does not involve \(p\), so neither does the limit.  Finally the
sentences of limit \(1\) are exactly those with \(T\vdash B\): the deductive
closure of \(T\).
\end{proof}
\end{annexitem}

\begin{annexitem}{annex:p03-w03-p03-067-categorical}{thm:p03-w03-p03-067-categorical}{The constant-\(p\) theory is \(\aleph_0\)-categorical}
\begin{proof}
\emph{No finite model.}  A structure has a nonempty universe.  Let \(H\) be
a graph with exactly \(m\geq1\) vertices and list them as
\(x_1,\ldots,x_m\).  They are distinct, so \(A_{m,0}\) requires a vertex
distinct from all of them, and there is none.  So \(H\) fails \(A_{m,0}\).

\emph{A countable model.}  By \cref{thm:p03-w03-p03-067-zero-one}, for any
constant \(p\in(0,1)\) every graph sentence holds almost surely or almost never
in \(G(n,p)\), and the almost sure theory is the deductive closure of \(T\).
By item~3 of \cref{thm:p03-w03-p03-067-as-theory} it has a countably infinite
model, which is a countable model of \(T\).

\emph{One-point extension.}  Let \(G,G'\) be models of \(T\) and \(\Phi\) a
finite partial isomorphism from \(G\) to \(G'\): an injective map from a
finite set of vertices of \(G\) into \(G'\) with \(a\sim b\) iff
\(\Phi(a)\sim\Phi(b)\) on its domain.  Let \(x\in G\) lie outside the domain.
Put \(P=\{\Phi(a):a\sim x\}\) and \(Q=\{\Phi(a):a\not\sim x\}\), \(a\)
ranging over the domain.  They partition the range of \(\Phi\), whose points
are distinct.  By \(A_{\abs P,\abs Q}\) in \(G'\) there is \(x'\) outside the
range, adjacent to every point of \(P\) and to none of \(Q\).  Then
\(\Phi\cup\{(x,x')\}\) is injective, and for each \(a\) in the domain,
\(a\sim x\) iff \(\Phi(a)\in P\) iff \(\Phi(a)\sim x'\); with irreflexivity
and symmetry it is again a partial isomorphism.  Applying the same argument
to \(\Phi^{-1}\), with \(A_{r,s}\) in \(G\), adds any given point of \(G'\)
outside the range to the range.

\emph{Back and forth.}  Let \(G,G'\) be countable models of \(T\).  Both are
infinite, so list their vertices without repetition as
\(g_1,g_2,\ldots\) and \(g'_1,g'_2,\ldots\).  Put \(\Phi_0=\varnothing\).  At
the Left Stage \(2k-1\), let \(x\) be the vertex \(g_i\) of least index outside
the domain of \(\Phi_{2k-2}\) and extend to \(\Phi_{2k-1}\) with \(x\) in its
domain.  At the Right Stage \(2k\), let \(x'\) be the vertex \(g'_j\) of least
index outside the range of \(\Phi_{2k-1}\) and extend to \(\Phi_{2k}\) with
\(x'\) in its range.  By induction on \(k\), \(g_1,\ldots,g_k\) lie in the
domain and \(g'_1,\ldots,g'_k\) in the range of \(\Phi_{2k}\): if
\(g_1,\ldots,g_{k-1}\) are already in the domain, then either \(g_k\) is
too, or it is the least missing vertex and is added at stage \(2k-1\); the
range is handled at stage \(2k\) in the same way.  The union \(\Phi\) of the
increasing chain is a function with domain all of \(G\) and range all of
\(G'\).  Any two vertices lie in the domain of a single \(\Phi_k\), so
\(\Phi\) is injective and preserves adjacency and nonadjacency.  It is an
isomorphism.

\emph{Identification.}  For a simple graph, the extension property of
\cref{thm:p03-random-graph}, for disjoint finite \(U\) and \(V\), is exactly
the axiom \(A_{\abs U,\abs V}\) applied to listings of \(U\) and \(V\).  So a
countable simple graph has that property iff it is a model of \(T\), and the
countable model of \(T\), unique by the previous paragraph, is the graph
characterized there.
\end{proof}
\end{annexitem}

\begin{annexitem}{annex:p03-w03-p03-067-rado-limit}{thm:p03-w03-p03-067-rado-limit}{Limits are truth values in the Rado graph}
\begin{proof}
\(R\) is a model of \(T\) by \cref{thm:p03-w03-p03-067-categorical}.  By
\cref{thm:p03-w03-p03-067-complete}, \(T\vdash B\) or \(T\vdash\neg B\).  If
\(R\entails B\) then \(T\vdash\neg B\) is impossible, by soundness, so
\(T\vdash B\) and the limit is \(1\) by
\cref{thm:p03-w03-p03-067-zero-one}.  If \(R\entails\neg B\) then
\(T\vdash\neg B\) and the limit is \(0\).
\end{proof}
\end{annexitem}

\begin{annexitem}{annex:p03-w03-p03-067-unary}{thm:p03-w03-p03-067-unary}{The unary regime and its countable model}
\begin{proof}
\emph{Part 1.}  Fix \(r\geq1\) and let \(X\) be the number of points
satisfying \(U\) in \(SU(n,p(n))\).  It is a sum of \(n\) independent
indicators of probability \(p=p(n)\), so \(\mathbb E X=np\) and
\(\operatorname{Var}X=np(1-p)\leq np\).  The sentence \(U_r\) fails exactly
when \(X\leq r-1\).  Put \(\mu=np\).  Since \(\mu\to\infty\), eventually
\(\mu>r-1\), and then \(X\leq r-1\) implies \(\abs{X-\mu}\geq\mu-r+1>0\).
Chebyshev's inequality gives
\[
 \Pr[X\leq r-1]\leq\frac{\mu}{(\mu-r+1)^2}\longrightarrow0
 \qquad(\mu\to\infty).
\]
The number \(n-X\) of points satisfying \(\neg U\) is a sum of \(n\)
independent indicators of probability \(1-p\), and the same estimate with
\(n(1-p)\to\infty\) shows that \(N_r\) holds almost surely.  A constant
\(p\in(0,1)\) gives \(np\to\infty\) and \(n(1-p)\to\infty\).

\emph{Part 2.}  A finite subset of \(T^{SU}\) is contained in
\(\{U_r,N_r:r\leq m\}\) for some \(m\), and a set of \(2m\) points with exactly
\(m\) of them in \(U\) satisfies it; by \cref{thm:fol-compactness}
\(T^{SU}\) has a model.  In any model \(\mathcal M\), the sentences \(U_r\)
and \(N_r\) for all \(r\) make \(U^{\mathcal M}\) and its complement
infinite.  Let \(\mathcal M,\mathcal M'\) be models and \(t\geq1\).  In the
\(t\)-round game of \cref{def:p03-ef-game} for the language \(\{U\}\),
Duplicator answers a repeated choice with its earlier partner, and a new
choice \(a\) with a vertex of the other structure that has not been chosen and
has the same \(U\)-status as \(a\).  One exists, because only finitely many
points have been chosen and each status class is infinite.  The pairing
then preserves equality and \(U\), so Duplicator wins.  By \cref{thm:p03-ef},
any two models agree on all sentences, and as in the proof of
\cref{thm:p03-w03-p03-067-complete}, \(T^{SU}\vdash A\) or
\(T^{SU}\vdash\neg A\) for every sentence \(A\) by
\cref{thm:fol-completeness}.  If \(T^{SU}\vdash A\), finitely many sentences
of \(T^{SU}\) entail \(A\), and part~1 with the union bound gives
\(\Pr[SU(n,p(n))\entails A]\to1\); if \(T^{SU}\vdash\neg A\), the limit is
\(0\).  The value is fixed by whether \(T^{SU}\vdash A\), which does not
involve \(p(n)\).

\emph{Part 3.}  A countable structure for \(\{U\}\) satisfies every \(U_r\)
and \(N_r\) iff \(U\) and its complement are both infinite: the forward
direction was shown in part~2, and the converse is immediate.  Let
\(\mathcal M,\mathcal M'\) be countable models.  Then \(U^{\mathcal M}\) and
\(U^{\mathcal M'}\) are countably infinite, so there is a bijection \(f\)
between them, and the complements are countably infinite, so there is a
bijection \(g\) between them.  The union \(f\cup g\) is a bijection from
\(\mathcal M\) onto \(\mathcal M'\) that maps \(U\)-points to \(U\)-points and
the other points to other points, so it is an isomorphism.  A countable model
exists: the natural numbers with \(U\) the even numbers.
\end{proof}
\end{annexitem}

\begin{annexitem}{annex:p03-w03-p03-067-sparse-failure}{thm:p03-w03-p03-067-sparse-failure}{Polynomially small \(p\) defeats an extension axiom}
\begin{proof}
Take \(n>r\) so large that \(p(n)\leq n^{-\alpha}\), and fix \(r\) distinct
vertices \(v_1,\ldots,v_r\).  If the outcome satisfies \(A_{r,0}\), some
\(z\notin\{v_1,\ldots,v_r\}\) is adjacent to every \(v_i\).  For a fixed such
\(z\) the \(r\) pairs \(\{z,v_i\}\) are distinct, so this has probability
\(p^r\).  By the union bound over the \(n-r\) candidates,
\[
 \Pr\bigl[G(n,p(n))\entails A_{r,0}\bigr]\leq(n-r)p(n)^r\leq n\cdot n^{-r\alpha}
 =n^{1-r\alpha}\longrightarrow0,
\]
since \(r\alpha>1\).  For \(p=n^{-\alpha}\) with \(\alpha>1/3\), the integer
\(r=3\) satisfies \(3>1/\alpha\).
\end{proof}
\end{annexitem}

\begin{annexitem}{annex:p03-w03-p03-067-axioms-generate}{thm:p03-w03-p03-067-axioms-generate}{The two schemas axiomatize \(T_\alpha\)}
\begin{proof}
\emph{Rigidity in the two presentations.}  Let \(X'\) and \(Z\) be disjoint
finite sets of vertices of a graph \(G\), with \(Z\) nonempty.  For a set
\(Y\) disjoint from \(X\), the count \(e_G(X,Y)\) of
\cref{def:p03-w03-p03-070-weight} and the count \(e_G(Y;X)\) of
\cref{def:p03-w03-p03-068-standing} are both the number of edges of \(G\)
with an end in \(Y\) and both ends in \(X\cup Y\), so
\(w_\alpha(X,Y)=\abs Y-e_G(Y;X)\,\alpha\).  The maps \(S\mapsto S\setminus
X'\) and \(Z_0\mapsto X'\cup Z_0\) are mutually inverse bijections between
the sets \(S\) with \(X'\subseteq S\subsetneq X'\cup Z\) and the sets
\(Z_0\subsetneq Z\), and under them \(Z\setminus S=Z\setminus Z_0\).  So the
condition \(\abs{Z\setminus S}-e_G(Z\setminus S;S)\,\alpha<0\) for all such
\(S\), which says that \(Z\) is rigid over \(X'\) in the sense of
\cref{def:p03-w03-p03-068-rigid-over}, is the condition
\(w_\alpha(X'\cup Z_0,Z\setminus Z_0)<0\) for all \(Z_0\subsetneq Z\), which
says that \(Z\) is \(\alpha\)-rigid over \(X'\) in the sense of
\cref{def:p03-w03-p03-070-weight}.

\emph{Item 1.}  Fix a graph \(G\).  For a finite graph \(H\) with
\(v-e\alpha<0\), the Nonexistence axiom of
\cref{def:p03-w03-p03-068-axioms}, with copies read as injective
edge-preserving maps, is word for word the Nonexistence axiom of
\cref{def:p03-w03-p03-070-theory} for \(F=H\).

For Generic Extension, the data agree: by item~2 of
\cref{thm:p03-w03-p03-070-same-theory}, a rooted graph is safe in the sense
of \cref{def:p03-w03-p03-065-dense-sparse} if and only if it is
\(\alpha\)-safe.  Fix a safe \((R,H)\) with roots \(x_1,\ldots,x_r\) and
nonroots \(y_1,\ldots,y_v\), an integer \(t\geq0\), distinct vertices
\(a_1,\ldots,a_r\) of \(G\), and distinct vertices \(b_1,\ldots,b_v\) of
\(G\) distinct from the \(a_i\).  Compare the three requirements on
\(\bar b\) over \(\bar a\).
\begin{itemize}
\item The \(b_j\) form an \((R,H)\) extension over the \(a_i\)
  (\cref{def:p03-w03-p03-068-rigid-over}) exactly when \(b_i\sim b_j\) for
  \(y_iy_j\in E(H)\) and \(a_i\sim b_j\) for \(x_iy_j\in E(H)\); this is
  clause~(1) of \cref{def:p03-w03-p03-070-theory}, every edge of \(H\)
  meeting a nonroot being present.
\item The absence of adjacencies \(b_ib_j\) and \(a_ib_j\) other than those of
  \(H\) is clause~(2).
\item Let \(t\geq1\), \(A=\{a_1,\ldots,a_r\}\) and \(B=\{b_1,\ldots,b_v\}\).
  Genericity in the sense of \cref{def:p03-w03-p03-068-axioms} quantifies
  over pairwise distinct \(z_1,\ldots,z_w\) outside \(A\cup B\) with
  \(w\leq t\) and \(\{z_1,\ldots,z_w\}\) rigid over \(A\cup B\); rigidity
  over a set is defined only for nonempty sets, so \(w\geq1\).  The sets of
  entries of such tuples are exactly the nonempty sets \(Z\) of at most
  \(t\) vertices outside \(A\cup B\) that are rigid over \(A\cup B\).  By the
  first paragraph these are the nonempty sets of at most \(t\) vertices
  outside \(A\cup B\) that are \(\alpha\)-rigid over \(A\cup B\); the empty
  set contains no vertex, so it never violates clause~(3).  Both definitions
  then forbid an edge from a vertex of such a \(Z\) to a vertex of \(B\).  So
  the two genericity conditions coincide.
\end{itemize}
For \(t=0\) neither axiom imposes genericity.  Hence for each choice of
\(\bar a\) the two axioms ask for the existence of the same tuples
\(\bar b\), and \(G\) satisfies one axiom if and only if it satisfies the
other.  Both \(\Sigma_\alpha\) and \(T^-_\alpha\) contain the graph axioms,
so their models are graphs, and among graphs they are satisfied together
axiom by axiom.  So they have the same models.

\emph{Item 2.}  Every outcome of \(G(n,n^{-\alpha})\) is a graph, so the graph
axioms hold with probability \(1\).  Let \(\varphi\) be the Nonexistence or
Generic Extension axiom of \cref{def:p03-w03-p03-068-axioms} for some datum,
and \(\varphi'\) the axiom of \cref{def:p03-w03-p03-070-theory} for the same
datum.  By item~1, \(\Pr[G(n,n^{-\alpha})\entails\varphi]=
\Pr[G(n,n^{-\alpha})\entails\varphi']\) for every \(n\), and the right-hand
side tends to \(1\) by \cref{thm:p03-w03-p03-070-as-holds}.

\emph{Item 3.}  By item~2 of \cref{thm:p03-w03-p03-070-complete}, for every
sentence \(A\) the limit \(f_A(\alpha)\) of
\(\Pr[G(n,n^{-\alpha})\entails A]\) exists and equals \(1\) if
\(T^-_\alpha\derives A\) and \(0\) otherwise.  So \(A\) holds almost surely
if and only if \(T^-_\alpha\derives A\).  By soundness and completeness
(\cref{thm:fol-soundness,thm:fol-completeness}), \(T^-_\alpha\derives A\) if
and only if every model of \(T^-_\alpha\) satisfies \(A\), and by item~1 this
holds if and only if every model of \(\Sigma_\alpha\) satisfies \(A\), that
is, \(A\) is a logical consequence of \(\Sigma_\alpha\).  Hence \(T_\alpha\)
is the set of logical consequences of \(\Sigma_\alpha\).

The same item of \cref{thm:p03-w03-p03-070-complete} states that
\(T^-_\alpha\) is consistent and complete.  A model of \(T^-_\alpha\) is a
model of \(\Sigma_\alpha\) by item~1, and it satisfies every logical
consequence of \(\Sigma_\alpha\), so \(T_\alpha\) is consistent.  For each
sentence \(A\), completeness gives \(T^-_\alpha\derives A\) or
\(T^-_\alpha\derives\neg A\), so \(A\in T_\alpha\) or \(\neg A\in T_\alpha\),
and \(T_\alpha\) is complete.  Finally let \(M\) be a structure for
\(\{\sim\}\).  If \(M\) is a model of \(\Sigma_\alpha\), it satisfies every
logical consequence of \(\Sigma_\alpha\), hence every sentence of
\(T_\alpha\).  Conversely \(\Sigma_\alpha\subseteq T_\alpha\) by item~2, so
a model of \(T_\alpha\) is a model of \(\Sigma_\alpha\), that is, a graph
satisfying every axiom of \cref{def:p03-w03-p03-068-axioms}.
\end{proof}
\end{annexitem}

\begin{annexitem}{annex:p03-w03-p03-067-diophantine}{thm:p03-w03-p03-067-diophantine}{Ceilings just above an irrational multiple}
\begin{proof}
For every integer \(u\geq1\) the number \(uc\) is not an integer, because
\(c\) is irrational; so \(\lceil uc\rceil-uc\in(0,1)\).  Put
\[
 \eta'=\min\bigl(\{\eta,\tfrac12\}\cup\{\lceil uc\rceil-uc:1\leq u<M\}\bigr)>0.
\]
It suffices to find \(v\geq1\) with \(\lceil vc\rceil-vc<\eta'\): such a
\(v\) is not any \(u<M\), since those have \(\lceil uc\rceil-uc\geq\eta'\),
so \(v\geq M\).

Let \(L\) be an integer with \(L\geq1/\eta'\), and write \(\{y\}=y-\lfloor
y\rfloor\).  The \(L+1\) numbers \(\{jc\}\), \(0\leq j\leq L\), lie in
\([0,1)\), which is the union of the \(L\) intervals
\([l/L,(l+1)/L)\); so two of them, say for \(i<j\), lie in one interval and
\(\delta=\{jc\}-\{ic\}\) satisfies \(\abs\delta<1/L\leq\eta'\).  Put
\(q=j-i\geq1\).  Then \(qc=(\lfloor jc\rfloor-\lfloor ic\rfloor)+\delta\),
and \(\delta\neq0\) because \(qc\) is not an integer.

If \(\delta<0\), then \(qc\) is an integer plus \(\delta\in(-1,0)\), so
\(\lceil qc\rceil-qc=-\delta\in(0,\eta')\), and \(v=q\) works.

If \(\delta>0\), let \(h=\lfloor(1-\eta')/\delta\rfloor+1\), the least
integer with \(h\delta>1-\eta'\).  Then
\(h\delta\leq1-\eta'+\delta<1\), since \(\delta<\eta'\).  Now
\(hqc=h(\lfloor jc\rfloor-\lfloor ic\rfloor)+h\delta\) is an integer plus
\(h\delta\in(1-\eta',1)\), so
\(\lceil hqc\rceil-hqc=1-h\delta\in(0,\eta')\), and \(v=hq\) works.
\end{proof}
\end{annexitem}

\begin{annexitem}{annex:p03-w03-p03-067-circulant}{thm:p03-w03-p03-067-circulant}{Rigid one-root extensions of nearly zero excess}
\begin{proof}
Write \(K=K(v)\), \(Q=Q_v\) and \(d=1/\alpha=k+c\).

\emph{The definition is sound.}  Since \(v\geq2k+1\), for each nonroot \(j\)
the \(2k\) nonroots at cyclic distance \(1,\ldots,k\) from \(j\) are
distinct and differ from \(j\); so \(K\) is a \(2k\)-regular simple graph
and has \(2kv/2=kv\) edges.  Since \(0<c<1\), each difference
\(\lceil jc\rceil-\lceil(j-1)c\rceil\) is \(0\) or \(1\), and summing over
\(j=1,\ldots,v\) gives \(\abs Q=\lceil vc\rceil-\lceil0\rceil=\lceil
vc\rceil\).  The edges of \(H(v)\) are those of \(K\) and the \(\abs Q\) root
edges, so \(e=kv+\lceil vc\rceil\).

\emph{Item 1.}  \(v-e\alpha=\alpha(dv-kv-\lceil vc\rceil)=\alpha(cv-\lceil
vc\rceil)=-\alpha g(v)\), and \(g(v)>0\) because \(vc\) is not an integer.

\emph{Counting \(Q\) on arcs.}  Call a set of nonroots an \emph{arc} if it is
nonempty, is not all of \(\{1,\ldots,v\}\), and consists of consecutive
nonroots in the cyclic order \(1,2,\ldots,v,1\).  For an arc \(I\),
\[
 \abs{Q\cap I}>c\abs I-1.
\]
First let \(I=\{i+1,\ldots,i+L\}\) with \(0\leq i<i+L\leq v\).  Summing the
differences gives \(\abs{Q\cap I}=\lceil(i+L)c\rceil-\lceil ic\rceil\), and
\(\lceil(i+L)c\rceil\geq(i+L)c\) while \(\lceil ic\rceil<ic+1\); so
\(\abs{Q\cap I}>Lc-1\).  The same two facts the other way round give
\(\abs{Q\cap I}<Lc+1\) for such \(I\).  Otherwise \(I\) contains \(v\) and
\(1\), and its complement \(J\) is nonempty, contains neither, and is of the
first form.  Then
\(\abs{Q\cap I}=\abs Q-\abs{Q\cap J}>\lceil vc\rceil-c\abs J-1\geq
c(v-\abs J)-1=c\abs I-1\).

\emph{Item 2.}  Let \(\{x\}\subseteq S\subsetneq V(H(v))\) and let \(W\) be
the nonempty set of nonroots outside \(S\).  The rooted graph \((S,H(v))\)
has \(\abs W\) nonroots, and its edge parameter \(e_W\) is the number of
edges of \(H(v)\) with an end in \(W\).  We show \(e_W>d\abs W\), which is
\(\abs W-e_W\alpha<0\).  Let \(\partial W\) be the number of edges of \(K\)
with exactly one end in \(W\).  Counting the degrees of the vertices of
\(W\) in \(K\) gives \(2k\abs W=2e(K|_W)+\partial W\).  The edges of
\(H(v)\) with an end in \(W\) are the edges of \(K|_W\), the \(\partial W\)
edges leaving \(W\), and the root edges to \(Q\cap W\), so
\[
 e_W=e(K|_W)+\partial W+\abs{Q\cap W}=k\abs W+\tfrac12\partial
 W+\abs{Q\cap W}.
\]
Since \(d=k+c\), it remains to show
\(\tfrac12\partial W+\abs{Q\cap W}>c\abs W\).  If \(W\) is all of the
nonroots, then \(\partial W=0\) and \(\abs Q=\lceil vc\rceil>vc\).
Otherwise \(W\) is the disjoint union of its maximal arcs
\(I_1,\ldots,I_m\), \(m\geq1\).  For each \(I_l\), the cycle edges joining
its first element to the preceding nonroot and its last element to the
following nonroot are edges of \(K\), since \(k\geq1\); they are two
different edges, since \(v\geq3\); and their other ends lie outside \(W\)
by maximality.  Each of these edges has exactly one end in \(W\), and that end
lies in \(I_l\), so the \(2m\) edges so obtained are distinct and
\(\partial W\geq2m\).  By the count on arcs,
\[
 \tfrac12\partial W+\abs{Q\cap W}\geq
 m+\sum_{l\leq m}\abs{Q\cap I_l}>m+\sum_{l\leq m}(c\abs{I_l}-1)=c\abs W.
\]

\emph{Item 3.}  Every vertex of \(K|_S\) has degree at most \(2k\), so
\(e(K|_S)\leq k\abs S<d\abs S\), using \(k<d\).  Multiplying by \(\alpha\)
gives \(e(K|_S)\alpha<\abs S\).

\emph{Item 4.}  Write \(N_i\) for the nonroots of the \(i\)-th block, so
\(V(H^s)=\{x\}\cup N_1\cup\cdots\cup N_s\), and every edge of \(H^s\) lies
inside one block \(\{x\}\cup N_i\).  Let \(U\subseteq V(H^s)\) be nonempty
and \(U_i=U\cap N_i\); safety of \((\varnothing,H^s)\) asks that
\(\abs U-e(H^s|_U)\alpha>0\).  If \(x\notin U\), the edges of \(H^s|_U\) are
those of the \(K(v_i)|_{U_i}\), and by item~3
\(\abs U-e(H^s|_U)\alpha=\sum_{U_i\neq\varnothing}(\abs{U_i}-e(K(v_i)|_{U_i})\alpha)>0\).
If \(x\in U\), put
\(\varphi_i(U_i)=\abs{U_i}-e(H(v_i)|_{\{x\}\cup U_i})\alpha\), so that
\(\abs U-e(H^s|_U)\alpha=1+\sum_{i\leq s}\varphi_i(U_i)\).  For
\(U_i\neq N_i\), rigidity (item~2) applied to \(S=\{x\}\cup U_i\) says
\(\abs{N_i\setminus U_i}-(e_i-e(H(v_i)|_S))\alpha<0\), where \(e_i\) is the
number of edges of \(H(v_i)\); rearranged, \(\varphi_i(U_i)>v_i-e_i\alpha\).
For \(U_i=N_i\) equality holds.  By item~1,
\(\varphi_i(U_i)\geq-\alpha g(v_i)>-\alpha2^{-i}\), so
\[
 1+\sum_{i\leq s}\varphi_i(U_i)>1-\alpha\sum_{i\geq1}2^{-i}=1-\alpha>0.
\]
So \((\varnothing,H^s)\) is safe.  A copy of \(H^s\) in a graph is exactly an
\((\varnothing,H^s)\) extension, so by
\cref{thm:p03-w03-p03-065-safe-criterion} the extension statement
\(\operatorname{Ext}(\varnothing,H^s)\) holds almost surely in
\(G(n,n^{-\alpha})\).
\end{proof}
\end{annexitem}

\begin{annexitem}{annex:p03-w03-p03-067-infinite-closure}{thm:p03-w03-p03-067-infinite-closure}{A countable model with an infinite closure}
\begin{proof}
Choose \(v_1,v_2,\ldots\) by \cref{thm:p03-w03-p03-067-diophantine} with
\(M=2k+1\) and \(\eta=2^{-i}\) for \(v_i\), and form \(H^s\) as in
\cref{thm:p03-w03-p03-067-circulant}.  By item~4 there,
\(\operatorname{Ext}(\varnothing,H^s)\in T_\alpha\) for every \(s\).

Add a constant symbol \(c^\ast\) to the language of graphs.  For \(s\geq1\)
let \(\tau_s\) be the sentence saying that there are vertices, distinct from
each other and from \(c^\ast\), forming an \((\{x\},H^s)\) extension over
\(c^\ast\); that is, \(c^\ast\) is the image of \(x\) in a copy of \(H^s\).
Let \(\Gamma=T_\alpha\cup\{\tau_s:s\geq1\}\).

\emph{\(\Gamma\) is finitely satisfiable.}  By
\cref{thm:p03-w03-p03-065-shelah-spencer}, every sentence holds almost surely
or almost never in \(G(n,n^{-\alpha})\), so by item~1 of
\cref{thm:p03-w03-p03-067-as-theory} \(T_\alpha\) has a model \(G\).  Let
\(\Gamma_0\subseteq\Gamma\) be finite and let \(s^\ast\) be the largest
\(s\) with \(\tau_s\in\Gamma_0\) (or \(s^\ast=1\) if there is none).  Since
\(G\) satisfies \(\operatorname{Ext}(\varnothing,H^{s^\ast})\), it contains a
copy of \(H^{s^\ast}\); interpret \(c^\ast\) as the image of \(x\).  For
\(s\leq s^\ast\), the graph \(H^s\) is the subgraph of \(H^{s^\ast}\)
formed by its first \(s\) blocks, and restricting the copy to these blocks
gives a copy of \(H^s\) with the same image of \(x\).  So the expanded
structure satisfies \(\Gamma_0\).

\emph{A countable model.}  By clause~1 of \cref{thm:fol-compactness},
\(\Gamma\) has a model \((G',c')\).  The language is finite, so
\cref{thm:downward-ls} with \(A=\{c'\}\) gives a countable elementary
substructure \((G_1,c')\) of \((G',c')\) containing \(c'\).  It satisfies the
sentences of \(\Gamma\), so its graph reduct \(G_1\) is a model of
\(T_\alpha\), infinite by item~2 of \cref{thm:p03-w03-p03-067-as-theory},
hence countably infinite, and \(c'\) is the image of \(x\) in a copy of
\(H^s\) in \(G_1\) for every \(s\).

\emph{\(\mathrm{cl}(\{c'\})\) is infinite.}  Fix \(s\) and such a copy, and
let \(Y_i\), \(i\leq s\), be the image of the nonroots of the \(i\)-th block.
The sets \(Y_1,\ldots,Y_s\) are nonempty, pairwise disjoint and do not
contain \(c'\), since the copy is injective.  Each \(Y_i\) is rigid over
\(\{c'\}\) in \(G_1\) in the sense of \cref{def:p03-w03-p03-068-rigid-over}:
let \(\{c'\}\subseteq S\subsetneq\{c'\}\cup Y_i\) and let \(S_0\) be the
corresponding set of vertices of \(H(v_i)\).  The copy maps the edges of
\(H(v_i)\) with an end outside \(S_0\) to distinct edges of \(G_1\) with an
end in \(Y_i\setminus S\) and both ends in \(\{c'\}\cup Y_i\), so
\(e_{G_1}(Y_i\setminus S;S)\) is at least the edge parameter of
\((S_0,H(v_i))\), while \(\abs{Y_i\setminus S}\) is its number of nonroots.
Rigidity of \((\{x\},H(v_i))\) (item~2 of
\cref{thm:p03-w03-p03-067-circulant}) then gives
\(\abs{Y_i\setminus S}-e_{G_1}(Y_i\setminus S;S)\alpha<0\).  So
\(\{c'\}\subsetneq\{c'\}\cup Y_i\) is a rigid \(v_i\)-chain from
\(\{c'\}\), and by item~2 of \cref{thm:p03-w03-p03-068-well-defined},
\(Y_i\subseteq\mathrm{cl}_{v_i}(\{c'\})\subseteq\mathrm{cl}(\{c'\})\).
Hence \(\mathrm{cl}(\{c'\})\) has at least \(1+s\) elements.  As \(s\) is
arbitrary, it is infinite.
\end{proof}
\end{annexitem}

\begin{annexitem}{annex:p03-w03-p03-067-demand-model}{thm:p03-w03-p03-067-demand-model}{The witness-demand graph models \(T_\alpha\) and has finite closures}
\begin{proof}
Write \(e(Y;S)\) for \(e_{G^\ast}(Y;S)\).  The vertex set of \(G^\ast\) is
\(\bigcup_mP_m\), and each vertex \(u\) lies in exactly one fresh set
\(F_{j(u)}\).

\emph{Where the edges are.}  The edges added at stage \(l\) join a vertex of
\(F_l\) to a vertex of \(F_l\) or to a root of the demand met at stage
\(l\), and these roots lie in \(P_{l-1}\).  Let \(u\in F_l\).  An edge at
\(u\) was added at stage \(l\), or at a later stage \(l'\) at which \(u\) was
a root, and then its other end lies in \(F_{l'}\), outside \(P_l\).  Hence,
for every \(l\geq1\): (b) an edge of \(G^\ast\) at a vertex of \(F_l\) whose
other end lies in \(P_l\) was added at stage \(l\), so it is the image of an
edge of the demand graph met at stage \(l\) and its other end lies in
\(F_l\) or among that demand's roots; (c) every other edge at a vertex of
\(F_l\) has its other end outside \(P_l\).

\emph{Item 1.}  Induction on \(\abs Y\).  Let \(j\) be the largest
\(j(u)\) for \(u\in Y\); then \(j>m\), since \(Y\) avoids \(P_m\).  Put
\(Y_j=Y\cap F_j\neq\varnothing\) and \(Y'=Y\setminus Y_j\).  The edges
counted by \(e(Y;P_m)\) are those within \(P_m\cup Y\) with an end in
\(Y_j\), and those within \(P_m\cup Y'\) with an end in \(Y'\), counted by
\(e(Y';P_m)\).  An edge of the first kind has both ends in
\(P_m\cup Y\subseteq P_j\), so by (b) it was added at stage \(j\) and is the
image of an edge of the stage-\(j\) demand graph \((R_j,H_j)\) with an end
at a nonroot whose image lies in \(Y_j\); its other end is a root or a
nonroot whose image lies in \((P_m\cup Y)\cap F_j=Y_j\).  Distinct edges have
distinct preimages.  Let \(L\) be the set of nonroots of \(H_j\) whose
images lie in \(Y_j\), so \(\abs L=\abs{Y_j}\).  The number of edges of the
first kind is therefore at most the edge parameter \(e_0\) of the
subextension \((R_j,H_j|_{R_j\cup L})\), which is sparse because
\((R_j,H_j)\) is safe; so \(\abs{Y_j}-e_0\alpha>0\).  If
\(Y'=\varnothing\) we are done; otherwise
\(\abs{Y'}-e(Y';P_m)\alpha>0\) by induction, and adding the two
inequalities gives \(\abs Y-e(Y;P_m)\alpha>0\).

\emph{Item 2.}  Each stage adds at least one vertex, since every rooted graph
has \(v\geq1\).  A stage always has a demand to meet: for each \(m\geq1\)
the edgeless graph on \(m\) vertices with no roots is safe, since each of
its nonempty subsets has parameter equal to its size, and these
infinitely many demands have no roots.  So the run has infinitely many
stages and \(G^\ast\) is countably infinite.  Let a demand \(\Delta\) have
roots in \(\bigcup_mP_m\); they lie in some \(P_{m_0}\).  Let \(\Delta\) be
the \(N\)-th demand of the enumeration.  At every stage after \(m_0\) at
which \(\Delta\) is not yet met, the demand met is unmet, eligible and
earlier than \(\Delta\) in the enumeration; there are fewer than \(N\) such
demands and each is met once, so \(\Delta\) is met within \(N\) stages
after \(m_0\).

\emph{Item 3.}  Let \(((R,H),\bar a)\) be met at stage \(m\) with fresh set
\(F_m=\{b_1,\ldots,b_v\}\).  The \(a_i\) lie in \(P_{m-1}\) and the \(b_j\)
in \(F_m\), so by (b) with \(l=m\) every edge of \(G^\ast\) among the \(b_j\) or
between a \(b_j\) and an \(a_i\) was added at stage \(m\) and is required by
\(H\).  Now let \(t\geq1\) and let \(Z\) be a set of at most \(t\) vertices,
disjoint from \(X=\{a_1,\ldots,a_r\}\) and from \(F_m\), rigid over
\(X\cup F_m\), and suppose some \(z\in Z\) is adjacent to some \(b_j\).  By
(b) and (c) with \(l=m\), \(z\in X\cup F_m\) or \(z\notin P_m\); the first is excluded, so
\(z\notin P_m\).  So \(Z'=Z\setminus P_m\) is nonempty, and by
\cref{thm:p03-w03-p03-068-move-base} with the base \(P_m\supseteq
X\cup F_m\), \(Z'\) is rigid over \(P_m\).  Taking \(S=P_m\) in the
definition of rigidity gives \(\abs{Z'}-e(Z';P_m)\alpha<0\), contrary to
item~1.  So no such \(Z\) exists, and \(b_1,\ldots,b_v\) are \(t\)-generic
over \(a_1,\ldots,a_r\).

\emph{Item 4.}  Nonexistence: let \(H\) be a finite graph with \(v\)
vertices, \(e\) edges and \(v-e\alpha<0\), and suppose \(G^\ast\) contains a
copy of \(H\) on a vertex set \(U\).  Then \(\abs U=v\) and the copy gives
\(e\) distinct edges within \(U\), so
\(\abs U-e(U;\varnothing)\alpha\leq v-e\alpha<0\), contrary to item~1 with
\(m=0\).  Generic Extension: let \((R,H)\) be safe, \(t\geq0\), and
\(a_1,\ldots,a_r\) distinct vertices of \(G^\ast\).  Replace \((R,H)\) by its
representative, relabelling the roots so that \(a_i\) corresponds to
\(x_i\).  By item~2 the demand is met, and by item~3 its fresh vertices
witness the axiom.  So \(G^\ast\) satisfies every axiom of
\cref{def:p03-w03-p03-068-axioms}.  It is a graph, so by item~3 of
\cref{thm:p03-w03-p03-067-axioms-generate} it is a model of
\(T_\alpha\).

\emph{Item 5.}  Let \(X\subseteq P_m\) and \(t\geq0\).  If a set \(Y\) with
\(1\leq\abs Y\leq t\) and \(Y\cap P_m=\varnothing\) were rigid over \(P_m\),
then taking \(S=P_m\) would give \(\abs Y-e(Y;P_m)\alpha<0\), contrary to
item~1.  So \(P_m\) is \(t\)-closed, and since \(\mathrm{cl}_t(X)\) is the
least \(t\)-closed set containing \(X\)
(\cref{thm:p03-w03-p03-068-well-defined}), \(\mathrm{cl}_t(X)\subseteq P_m\).
This holds for every \(t\), so \(\mathrm{cl}(X)\subseteq P_m\), which is
finite.  Every finite \(X\) lies in some \(P_m\), because the \(P_m\)
increase and exhaust the vertex set.
\end{proof}
\end{annexitem}

\begin{annexitem}{annex:p03-w03-p03-067-demand-unique}{thm:p03-w03-p03-067-demand-unique}{The witness-demand graph is unique}
\begin{proof}
For a graph \(G\), say that a finite set \(A\) of vertices satisfies
\((\ast)\) in \(G\) if \(\abs Y-e_G(Y;A)\alpha>0\) for every finite nonempty
set \(Y\) of vertices of \(G\) outside \(A\).  Let \(G^\ast\) and \(G^{\ast\ast}\)
be the graphs of two runs, with stage sets \(P^\ast_m\) and
\(P^{\ast\ast}_m\).  By item~1 of \cref{thm:p03-w03-p03-067-demand-model},
every \(P^\ast_m\) satisfies \((\ast)\) in \(G^\ast\) and every
\(P^{\ast\ast}_m\) in \(G^{\ast\ast}\); in particular \(\varnothing=P_0\) does.

\emph{Step 1: a met demand over a \((\ast)\) set gives a \((\ast)\) set.}  Let
\(A'\) satisfy \((\ast)\) in \(G^{\ast\ast}\), and let a demand whose roots
are the elements of \(A'\) be met at stage \(m\) with fresh set \(D\).  We
show that \(A'\cup D\) satisfies \((\ast)\).  Let \(Y\) be finite, nonempty
and disjoint from \(A'\cup D\); put \(Y_0=Y\cap P^{\ast\ast}_m\) and
\(Y_1=Y\setminus P^{\ast\ast}_m\).  By (b) and (c), with \(l=m\), in the proof of
\cref{thm:p03-w03-p03-067-demand-model}, no edge joins a vertex of \(D\) to a
vertex of \(P^{\ast\ast}_m\setminus(A'\cup D)\), so no edge joins \(Y_0\) to
\(D\).  Splitting the edges within \(A'\cup D\cup Y\) with an end in \(Y\)
according to whether they meet \(Y_1\),
\[
 e(Y;A'\cup D)=e(Y_0;A')+e(Y_1;A'\cup D\cup Y_0),
\]
where the first term omits \(D\) because no edge joins \(Y_0\) to \(D\), and
both counts are in \(G^{\ast\ast}\).  The roots lie in
\(P^{\ast\ast}_{m-1}\), so \(A'\cup D\cup Y_0\subseteq P^{\ast\ast}_m\), and
\(e(Y_1;A'\cup D\cup Y_0)\leq e(Y_1;P^{\ast\ast}_m)\), since enlarging the
base only adds edges.  If \(Y_0\neq\varnothing\),
\(\abs{Y_0}-e(Y_0;A')\alpha>0\) by \((\ast)\) for \(A'\); if
\(Y_1\neq\varnothing\), \(\abs{Y_1}-e(Y_1;P^{\ast\ast}_m)\alpha>0\) by
item~1 of \cref{thm:p03-w03-p03-067-demand-model}.  At least one is
nonempty, and adding gives \(\abs Y-e(Y;A'\cup D)\alpha>0\).

\emph{Step 2: the forth step.}  Let \(f\) be an isomorphism from
\(G^\ast|_A\) onto \(G^{\ast\ast}|_{A'}\), where \(A\) satisfies \((\ast)\) in
\(G^\ast\) and \(A'\) in \(G^{\ast\ast}\), and let \(u\) be a vertex of
\(G^\ast\) not in \(A\).  Choose \(m\) with \(A\cup\{u\}\subseteq P^\ast_m\)
and put \(B=P^\ast_m\).  The rooted graph \((A,G^\ast|_B)\) is safe: for
\(A\subsetneq S\subseteq B\), the subextension \((A,G^\ast|_S)\) has the
nonroots \(N=S\setminus A\) and edge parameter \(e_{G^\ast}(N;A)\), and
\(\abs N-e_{G^\ast}(N;A)\alpha>0\) by \((\ast)\) for \(A\).  List \(A\) as
\(a_1,\ldots,a_r\) and \(B\setminus A\) as \(u_1,\ldots,u_v\), and let
\((R,H)\) be the representative of this rooted graph, with \(x_i\)
corresponding to \(a_i\) and \(y_j\) to \(u_j\).  By item~2 of
\cref{thm:p03-w03-p03-067-demand-model}, the demand
\(((R,H),(f(a_1),\ldots,f(a_r)))\) is met in the second run, with fresh
vertices \(b_1,\ldots,b_v\) say; put \(D=\{b_1,\ldots,b_v\}\) and extend
\(f\) by \(u_j\mapsto b_j\).  The extended map is an isomorphism from
\(G^\ast|_B\) onto \(G^{\ast\ast}|_{A'\cup D}\): on \(A\) it is \(f\); an
edge \(u_iu_j\) or \(a_iu_j\) of \(G^\ast\) is an edge \(y_iy_j\) or
\(x_iy_j\) of \(H\), so \(b_ib_j\) or \(f(a_i)b_j\) is an edge of
\(G^{\ast\ast}\); and by item~3 of
\cref{thm:p03-w03-p03-067-demand-model}, every edge of \(G^{\ast\ast}\) among
the \(b_j\) or between a \(b_j\) and an \(f(a_i)\) comes from an edge of
\(H\), hence from an edge of \(G^\ast|_B\).  Its domain \(B\) satisfies
\((\ast)\) in \(G^\ast\), and its range \(A'\cup D\) satisfies \((\ast)\) in
\(G^{\ast\ast}\) by Step~1.  The back step, adding a given vertex of
\(G^{\ast\ast}\) to the range, is the same with the two runs exchanged and
\(f^{-1}\) in place of \(f\).

\emph{Step 3: back and forth.}  Both graphs are countably infinite by item~2
of \cref{thm:p03-w03-p03-067-demand-model}; list their vertices as
\(g_1,g_2,\ldots\) and \(g'_1,g'_2,\ldots\).  Start with \(f_0=\varnothing\),
whose domain and range \(\varnothing\) satisfy \((\ast)\).  At step
\(2l-1\), if \(g_l\) is not in the domain, apply Step~2 to add it; at step
\(2l\), if \(g'_l\) is not in the range, apply the back step to add it.  The
maps increase, and after step \(2l\) the vertices \(g_1,\ldots,g_l\) are in
the domain and \(g'_1,\ldots,g'_l\) in the range.  The union of the chain is
a bijection from the vertices of \(G^\ast\) onto those of \(G^{\ast\ast}\),
and any two vertices lie in the domain of one \(f_l\), which preserves
adjacency and nonadjacency.  So it is an isomorphism.
\end{proof}
\end{annexitem}

\begin{annexitem}{annex:p03-w03-p03-067-sparse-models}{thm:p03-w03-p03-067-sparse-models}{The sparse theory is not \(\aleph_0\)-categorical}
\begin{proof}
\Cref{thm:p03-w03-p03-067-infinite-closure} gives a countable model \(G_1\)
of \(T_\alpha\) with a vertex whose closure is infinite.
\Cref{thm:p03-w03-p03-067-demand-model} gives a countable model \(G^\ast\) of
\(T_\alpha\), the graph of any run, in which every finite vertex set has
finite closure.  Rigidity of a set over another, \(t\)-closedness,
\(t\)-closures and closures are defined from adjacency alone, so a graph
isomorphism \(\pi\) satisfies \(\pi[\mathrm{cl}(X)]=\mathrm{cl}(\pi[X])\).
An isomorphism from \(G_1\) onto \(G^\ast\) would send a vertex with
infinite closure to a vertex with infinite closure, and \(G^\ast\) has none.
So \(T_\alpha\) has two nonisomorphic countable models and is not
\(\aleph_0\)-categorical.
\end{proof}
\end{annexitem}
